% Owned proposed B.144. All checkpoint29 sources remain frozen.
\subsection*{The exact one-family fixed source and its full stabilizer}

\statement{B.144}{Proved}{The actual one-family fixed source is a coset-permutation free-by-free group and an explicit centralizer extension; its marked fixed section is equivalent to one whole-stabilizer commutation condition}
\label{stmt:B144}
Use B.139's actual groups and graph
\[
 \begin{gathered}
 Q_f=C_2(C)*\mathbb Z(f)*C_2(\tau),\qquad S=F(m,v),\qquad
 Q=Q_f\rtimes S,\\
 H=\langle m,fv\rangle=\{h_s^{-1}s:s\in S\},\qquad
 \mathcal G^{(0)}=Q/H,\\
 \mathcal G^{(1)}=\bigl\{\{qH,qCH\},\{qH,q\tau H\}:q\in Q\bigr\}.
 \end{gathered}
\]
Projection $H\to S$ is an isomorphism, and
\[
 E_a=A(\mathcal G)\rtimes Q,
 \qquad\Gamma_a=A(\mathcal G)\rtimes Q^0
\]
are the actual one-family orbifold quotient and its ordinary subgroup.
Write $\delta=\operatorname{Inn}(\tau)$, $a^*=\tau a\tau$,
$A=aa^*$, and define the \emph{entire actual stabilizer}
\begin{equation}\label{eq:B144-stabilizer}
 L=H\cap S=\{s\in S:h_s=1\}.
\end{equation}
Then the actual discrete fixed-source image has the exact descriptions
\begin{equation}\label{eq:B144-source-structure}
 J_a=\Gamma_a^\delta
   =F(S/L)\rtimes S
   \cong S*_{L}\bigl(L\times\mathbb Z(A)\bigr).
\end{equation}
The free basis indexed by the left cosets $sL$ is
\begin{equation}\label{eq:B144-norm-basis}
 A_{sL}=a_{sH}a_{\tau sH}=sAs^{-1}.
\end{equation}
In particular the kernel of the free-source map
$F(A,m,v)\twoheadrightarrow J_a$ is exactly
\[
 \langle\!\langle[A,l]:l\in L\rangle\!\rangle.
\]

In the original $\Gamma_0$, let
$A_{\rm src}=s_1s_4^{-1}s_1^{-1}$ be the actual corrected source word
of B.130. The marked assignment
\begin{equation}\label{eq:B144-marked-section}
 A\longmapsto A_{\rm src},\qquad m\longmapsto m,
 \qquad v\longmapsto v
\end{equation}
extends to a homomorphism $J_a\to\Gamma_0^\delta$ if and only if
\begin{equation}\label{eq:B144-whole-gate}
 [A_{\rm src},l]=1\quad\text{for every }l\in L.
\end{equation}
If it extends, it is a discrete section of the actual fixed-subgroup
surjection. The equality $L=\langle m\rangle$ and condition
\eqref{eq:B144-whole-gate} are not asserted. Nor is an abstract
completion of \eqref{eq:B144-source-structure}, a compact fixed section,
or the original fixed-source equality inferred.

\begin{proof}
\emph{B.144.1. Actual marking, the ordered cocycle, and parity.}
For clarity the coarse actions underlying the statement are exactly
\[
 \begin{array}{c|ccc}
     &C&f&\tau\\ \hline
 \alpha_m&f^{-1}Cf&f^{-1}CfCf&\tau\\
 \alpha_v&f^{-1}Cf&\tau f\tau&\tau.
 \end{array}
\]
The marked peripheral cocycle is
\begin{equation}\label{eq:B144-cocycle}
 \alpha_s(a)=h_sah_s^{-1},\quad h_m=1,\quad h_v=f^{-1},
 \quad h_{st}=\alpha_s(h_t)h_s,
 \quad h_{s^{-1}}=\alpha_{s^{-1}}(h_s^{-1}).
\end{equation}
These are B.139's complete signed operators and ordered action, not
a commutative approximation. Thus $H$ is precisely the stabilizer of
the basis vertex $a_H$, with the indicated projection isomorphism.
The $S$-orbit of $H$ is indexed by $S/L$, since
$sH=s'H$ exactly when $s^{-1}s'\in H\cap S$.
In particular $\langle m\rangle\subseteq L$, but this does not
identify the whole intersection.

The original orientation sends $C,\tau$ to $1\in C_2$ and
$f,m,v$ to $0$. It is trivial on $H$, so it defines a well-defined
bipartition of $Q/H$. Every edge of $\mathcal G$ crosses the two
parts, because both edge labels are odd. In particular the even
induced subgraph, on $Q^0/H$, has no edges. This is a property of
the actual infinite graph; it does not impose a bipartiteness
hypothesis on finite graph heads elsewhere in the proof of B.139.

The selected fixed-edge classification of B.139 gives all
$\tau$-exchanged edges as
\[
 \{sH,\tau sH\}\qquad(s\in S).
\]
One can see the entire classification directly: if
$\tau rH=rtH$ with $t=C$ or $\tau$, then
$r^{-1}\tau r\in tH$. Projection to $S$, and the injectivity of
$H\to S$, force the $H$ factor to be trivial. The two cone factors
are not conjugate, so $t=\tau$, and
$r\in C_Q(\tau)=\langle\tau\rangle\times S$.
Likewise no vertex is fixed by $\tau$, since its hypothetical
conjugate in $H$ would project trivially to $S$.
Distinct cosets $sL$ give distinct even endpoints and hence distinct
unordered exchanged pairs.

\emph{B.144.2. Simultaneous freeness of every actual norm.}
Kill all odd graph generators. The defining graph presentation gives
a retraction
\[
 A(\mathcal G)\longrightarrow F(Q^0/H),
\]
because the even induced subgraph has no edges. The norm in
\eqref{eq:B144-norm-basis} maps to the single free generator
$a_{sH}$. A nontrivial reduced word in distinct norm symbols therefore
remains a nontrivial reduced word in distinct even free generators.
Thus all these norms are simultaneously a free basis of their subgroup.

B.139's full finite-support fixed-word argument proves that the same
norms generate every element of $A(\mathcal G)^\delta$.
Combining generation with this retraction gives
\[
 A(\mathcal G)^\delta=F(S/L).
\]
The action of $S$ is the literal left coset permutation. There is no
source-orbit cutoff, no substitution of individually primitive elements
for a simultaneous basis, and no finite-word inference about compact words.

\emph{B.144.3. The exact source presentation.}
B.139 gives $(Q^0)^\delta=S$ and the exact actual equality
$J_a=A(\mathcal G)^\delta\rtimes S$.
This proves the first description in \eqref{eq:B144-source-structure}.
Consider the group
\[
 P=\langle S,A\mid[A,l]=1\text{ for every }l\in L\rangle.
\]
There is a homomorphism $P\to J_a$ with the stated marked images.
Conversely, send the free coset-basis element indexed by $sL$ to
$sAs^{-1}\in P$. This is well-defined precisely because $A$
commutes with $L$, and it is $S$-equivariant. The two maps are inverse
on $S$ and on every norm generator. Thus $P\cong J_a$.
The defining presentation of the amalgam
$S*_{L}(L\times\mathbb Z(A))$ is the same $P$, proving the final
description and the asserted full free-source kernel. This does not
assert a new amalgam embedding in the original unfilled group.

In particular the stabilizer of $A$ under the actual $S$ action is
exactly $L$, because distinct cosets label distinct free basis elements.
The presentation is not abbreviated to a single commutator with $m$.

\emph{B.144.4. The exact marked fixed-lifting criterion.}
B.130 proves that $A_{\rm src}$ is fixed, maps to $aa^*$, and
satisfies $[A_{\rm src},m]=1$. The actual sources $m,v$ are fixed
and retain their free source subgroup. The presentation in part~3
therefore proves that \eqref{eq:B144-marked-section} extends exactly
when \eqref{eq:B144-whole-gate} holds. If it extends, composition
with the actual fixed quotient is the identity on $A,m,v$, hence
on the whole $J_a$. It is consequently an injective section.

This is the criterion for the particular marked source lifts.
An arbitrary fixed section could choose different fixed lifts, so
failure of this criterion, if later proved, would not alone exclude
every possible fixed section. Conversely fixed surjectivity does not
by itself establish the criterion. Whole-group nonsplitting from
B.142 is compatible with a section on a selected fixed subgroup.

\emph{B.144.5. The restricted retraction and the compact boundary.}
There is an explicit $S$-equivariant retraction on the restricted group
$A(\mathcal G)\rtimes S$: send a chosen even generator $a_{sH}$ to
$A_{sL}$, every other graph generator to $1$, and fix $S$.
Every defining graph edge has a killed odd endpoint, so all relations
are respected. The map fixes each norm and each source, and therefore
retracts onto the displayed copy of $J_a$.

No extension of this retraction to all of $\Gamma_a$ is claimed;
retained $Q_f$ elements can move vertices outside this $S$ orbit.
In particular the argument does not identify the actual compact closure
of $J_a$ with its abstract completion, does not furnish a compact
fixed section, and does not insert a completed amalgam normal form.
B.139's compact equality remains the equality with the actual closed
source image proved there.
\end{proof}

The complete whole-stabilizer proof and independent full-source reviews
are retained with the actual B.139 marking, rather than an enumerated
return set:
\begin{quote}\footnotesize\raggedright
\path{work_CE/one_family_fixed_source_exact_presentation29.txt},
\path{work_A/review_one_family_fixed_source_presentation29.md},
\path{work_DE/next_marginal_source/review_one_family_fixed_source_presentation29.md}.
\end{quote}
